\section{AIME II 2001}
        \begin{enumerate}
        	\item (\textit{AIME II 2001}) Let $N$ be the largest positive integer with the following property: reading from left to right, each pair of consecutive digits of $N$ forms a perfect square. What are the leftmost three digits of $N$?


 
 \textbf{Answer: }816



\textbf{Solution}

The two-digit perfect squares are $16, 25, 36, 49, 64, 81$. We try making a sequence starting with each one:

 
                \begin{itemize}
                    \item $16 - 64 - 49$. This terminates since none of them end in a $9$, giving us $1649$.\item $25$.\item $36 - 64 - 49$, $3649$.\item $49$.\item $64 - 49$, $649$.\item $81 - 16 - 64 - 49$, $81649$.
                \end{itemize}


                The largest is $81649$, so our answer is $\boxed{816}$.

 




	\item (\textit{AIME II 2001}) Each of the 2001 students at a high school studies either Spanish or French, and some study both. The number who study Spanish is between 80 percent and 85 percent of the school population, and the number who study French is between 30 percent and 40 percent. Let $m$ be the smallest number of students who could study both languages, and let $M$ be the largest number of students who could study both languages. Find $M-m$. 


 
 \textbf{Answer: }298



\textbf{Solution}

Let $S$ be the number of people who study Spanish, $F$ be the number of people who study French, and let $S \cap F$ be the number of students who study both. Then $\left\lceil 80\% \cdot 2001 \right\rceil = 1601 \le S \le \left\lfloor 85\% \cdot 2001 \right\rfloor = 1700$, and $\left\lceil 30\% \cdot 2001 \right\rceil = 601 \le F \le \left\lfloor 40\% \cdot 2001 \right\rfloor = 800$. By the \href{https://artofproblemsolving.com/wiki/index.php?title=Principle\_of\_Inclusion-Exclusion}{Principle of Inclusion-Exclusion},

 
 \[S+F- S \cap F = S \cup F = 2001\]

 For $m = S \cap F$ to be smallest, $S$ and $F$ must be minimized.

 
 \[1601 + 601 - m = 2001 \Longrightarrow m = 201\]

 For $M = S \cap F$ to be largest, $S$ and $F$ must be maximized.

 
 \[1700 + 800 - M = 2001 \Longrightarrow M = 499\]

 Therefore, the answer is $M - m = 499 - 201 = \boxed{298}$.

 




	\item (\textit{AIME II 2001}) Given that

 \begin{align*}x_{1}&=211,\\ x_{2}&=375,\\ x_{3}&=420,\\ x_{4}&=523,\ \text{and}\\ x_{n}&=x_{n-1}-x_{n-2}+x_{n-3}-x_{n-4}\ \text{when}\ n\geq5, \end{align*}

find the value of $x_{531}+x_{753}+x_{975}$.


 
 \textbf{Answer: }898



\textbf{Solution}

We find that $x_5 = 267$ by the recursive formula. Summing the \href{https://artofproblemsolving.com/wiki/index.php?title=Recursion}{recursions}

 
 \begin{align*} x_{n}&=x_{n-1}-x_{n-2}+x_{n-3}-x_{n-4} \\ x_{n-1}&=x_{n-2}-x_{n-3}+x_{n-4}-x_{n-5} \end{align*}


 yields $x_{n} = -x_{n-5}$. Thus $x_n = (-1)^k x_{n-5k}$. Since $531 = 106 \cdot 5 + 1,\ 753 = 150 \cdot 5 + 3,\ 975 = 194 \cdot 5 + 5$, it follows that

 
 \[x_{531} + x_{753} + x_{975} = (-1)^{106} x_1 + (-1)^{150} x_3 + (-1)^{194} x_5 = 211 + 420 + 267 = \boxed{898}.\] 

 





\textbf{Solution Variant}

The recursive formula suggests telescoping. Indeed, if we add $x_n$ and $x_{n-1}$, we have $x_n + x_{n-1} = (x_{n-1} - x_{n-2} + x_{n-3} - x_{n-4}) + (x_{n-2} - x_{n-3} + x_{n-4} - x_{n-5}) = x_{n-1} - x_{n-5}$. 

 Subtracting $x_{n-1}$ yields $x_n = -x_{n-5} \implies x_n = -(-(x_{n-10})) = x_{n-10}$. 

 Thus,

 
 \[x_{531} + x_{753} + x_{975} = x_1 + x_3 + x_5 = x_1 + x_3 + (x_4 - x_3 + x_2 - x_1) = x_2 + x_4 = 375 + 523 = \boxed{898}.\]

 Notice that we didn't need to use the values of $x_1$ or $x_3$ at all.

 


 \textbf{Non-Rigorous Solution} Calculate the first few terms:

 
 \[211,375,420,523,267,-211,-375,-420,-523,\dots\]

 At this point it is pretty clear that the sequence is periodic with period 10 (one may prove it quite easily like in solution 1) so our answer is obviously $211+420+267=\boxed{898}$

 




	\item (\textit{AIME II 2001}) Let $R = (8,6)$. The lines whose equations are $8y = 15x$ and $10y = 3x$ contain points $P$ and $Q$, respectively, such that $R$ is the midpoint of $\overline{PQ}$. The length of $PQ$ equals $\frac {m}{n}$, where $m$ and $n$ are relatively prime positive integers. Find $m + n$.


 
 \textbf{Answer: }067



\textbf{Solution}


                    \begin{figure}[H]
                        \centering
                        \includegraphics[width=0.5\textwidth]{Solution/AIME_II_2001_4_Solution_Figure_1.png}
                    \end{figure}
                    The coordinates of $P$ can be written as $\left(8a, 15a\right)$ and the coordinates of point $Q$ can be written as $\left(10b, 3b\right)$. By the midpoint formula, we have $8a+10b=16$ and $15a+3b=12$. Substituting $b=4-5a$ we derive $a = \frac{4}{7}$ and $b = \frac{8}{7}$. Thus $P$ is $\left(\frac{32}7, \frac{60}7\right)$, Q is $\left(\frac{80}7, \frac{24}7\right)$, and the coordinates form a 3-4-5 triangle dilated by $\frac{12}7$. Finally the distance $PQ$ must be $\frac{60}{7}$ so the answer is $\boxed{067}$.

 




	\item (\textit{AIME II 2001}) A set of positive numbers has the $triangle~property$ if it has three distinct elements that are the lengths of the sides of a triangle whose area is positive. Consider sets $\{4, 5, 6, \ldots, n\}$ of consecutive positive integers, all of whose ten-element subsets have the triangle property. What is the largest possible value of $n$?


 
 \textbf{Answer: }253



\textbf{Solution}

Out of all ten-element subsets with distinct elements that do not possess the triangle property, we want to find the one with the smallest maximum element. Call this subset $\mathcal{S}$. Without loss of generality, consider any $a, b, c \,\in \mathcal{S}$ with $a < b < c$. $\,\mathcal{S}$ does not possess the \href{https://artofproblemsolving.com/wiki/index.php?title=Triangle\_inequality}{triangle property}, so $c \geq a + b$. We use this property to build up $\mathcal{S}$ from the smallest possible $a$ and $b$:

 
 \[\mathcal{S} = \{\, 4,\, 5,\, 4+5, \,5+(4+5),\, \ldots\,\} = \{4, 5, 9, 14, 23, 37, 60, 97, 157, 254\}\]

 $\mathcal{S}$ is the "smallest" ten-element subset without the triangle property, and since the set $\{4, 5, 6, \ldots, 253\}$ is the largest set of consecutive integers that does not contain this subset, it is also the largest set of consecutive integers in which all ten-element subsets possess the triangle property. Thus, our answer is $n = \fbox{253}$.

 \textbf{Note} If we wanted to find this for a much larger number (say 2001), we could have noted that this is a "quasi-Fibonacci" sequence with initial terms $4,5$ and built up an explicit function to find the $nth$ term. (The latter part is generally pretty annoying).

 




	\item (\textit{AIME II 2001}) Square $ABCD$ is inscribed in a circle. Square $EFGH$ has vertices $E$ and $F$ on $\overline{CD}$ and vertices $G$ and $H$ on the circle. The ratio of the area of square $EFGH$ to the area of square $ABCD$ can be expressed as $\frac {m}{n}$ where $m$ and $n$ are relatively prime positive integers and $m < n$. Find $10n + m$.


 
 \textbf{Answer: }251



\textbf{Solution 1}

Begin by drawing another square congruent to $EFGH$ but on side $AB$ instead of side $CD$. Call this square $IJKL$, such that $I$ lies on $\overline{AB}$ and $\overline{IJ}$ lies on line $FG$. Let the side length of $EFGH$ and $IJKL$ be $s$ and the side length of $ABCD$ be $1$.

 
                    \begin{figure}[H]
                        \centering
                        \includegraphics[width=0.5\textwidth]{Solution/AIME_II_2001_6_Solution_Figure_1.png}
                    \end{figure}
                    By Power of a Point on point $F$, we find that 

 
 \[FC \cdot FD = FG \cdot FJ\]
 \[\frac{1-s}{2} \cdot \frac{1+s}{2} = s(1+s)\]

 Expanding and simplifying we get the quadratic $5s^2 +4s-1 = 0$. Factoring, we get $(5s-1)(s+1) = 0$. Therefore, $s= \frac{1}{5}$ and the ratio of the squares' areas is $\frac{1}{25}$. Our answer is therefore $10 \cdot 25 + 1 = \boxed{251}$.

 





\textbf{Solution 2}

Let $O$ be the center of the circle, $2a$ be the side length of square $ABCD$, and $2b$ be the side length of square $EFGH$. By symmetry, the horizontal and vertical displacements of $C$ from $O$ are both $\frac{2a}{2} = a$, so by the \href{https://artofproblemsolving.com/wiki/index.php?title=Pythagorean\_Theorem}{Pythagorean theorem}, the radius of the circle is $OC = \sqrt{a^2+a^2} = a\sqrt{2}$.

 
                    \begin{figure}[H]
                        \centering
                        \includegraphics[width=0.5\textwidth]{Solution/AIME_II_2001_6_Solution_Figure_2.png}
                    \end{figure}
                    Now let $I$ be the midpoint of $\overline{GH}$, giving $\overline{GI} = \frac{2b}{2} = b$, and let $J$ be the point where $\overline{OI}$ intersects $\overline{CD}$. We observe that since a diameter bisects a chord perpendicular to it, $\overline{GH}$ must be perpendicular to the diameter passing through $I$. This means that triangle $OGI$ has a right angle at $I$, and that $\overline{OJ}$ and $\overline{JI}$ are both parallel to $\overline{BC}$ and $\overline{FG}$. As the horizontal displacement of $C$ from $O$ is $a$ (from above), it follows that $\overline{OJ} = a$, and hence 
 \[\overline{OI} = \overline{OJ}+\overline{JI} = a+\overline{FG} = a+2b,\] so by the Pythagorean Theorem again,

 
 \begin{align*} \overline{OG}^2 = \overline{OI}^2 + \overline{GI}^2 &\iff \left(a\sqrt{2}\right)^2 = (a+2b)^2 + b^2 \\ &\iff 2a^2 = a^2+4ab+4b^2+b^2 \\ &\iff a^2-4ab-5b^2 = 0 \\ &\iff (a-5b)(a+b) = 0 \end{align*}


 Since lengths must be positive, we clearly cannot have $a+b = 0$, so the solution must instead be $a = 5b$. Since the ratio of the areas of two similar figures is the square of the ratio of their corresponding side lengths, we deduce that the required ratio is $\left(\frac{1}{5}\right)^2 = \frac{1}{25}$, and so the answer is $10 \cdot 25 + 1 = \boxed{251}$.

 Alternative Finish for Solution 1:

 After establishing that $a^2-4ab-5b^2 = 0$, another way to proceed is to observe that what we actually need to calculate is $\frac{b}{a}$. Accordingly, we divide both sides of this equation by $a^2$ (or equivalently, choose units such that the area of square $ABCD$ is $1$, without loss of generality), giving
 \[1-4\left(\frac{b}{a}\right)-5\left(\frac{b}{a}\right)^2 = 0,\]which is a quadratic in precisely the variable $\frac{b}{a}$. Thus, by solving it, we immediately obtain $\frac{b}{a} = \frac{1}{5}$ or $\frac{b}{a} = -1$, and as in Solution 1a, the negative root is obviously extraneous. Hence the required ratio of areas is $\left(\frac{1}{5}\right)^2 = \frac{1}{25}$, and so the answer is $10 \cdot 25 + 1 = \boxed{251}$.

 


 





\textbf{Solution 3 (coordinates)}

Let $A$ be the top-left vertex of square $ABCD$, and label the rest of the vertices in alphabetical order going clockwise from $A$. Let $D$ have coordinates $(0,0)$ and the side length of square $ABCD$ be $a$. Let $\overline{DF} = b$, the midpoint of $\overline{EF}$ be $J$, and $\overline{PQ}$ be the diameter of the circle that passes through $J$. Since a diameter bisects a chord perpendicular to it, we deduce that $\overline{CJ} = \overline{DJ}$, while $\overline{JE} = \overline{JF}$ by the definition of the midpoint, so $\overline{CE} = \overline{CJ}+\overline{JE} = \overline{DJ}+\overline{JF} = \overline{DF} = b$. It follows that 
 \[2b = \overline{CE}+\overline{DF} = \left(\overline{CF}+\overline{FE}\right)+\left(\overline{DE}+\overline{FE}\right) = \left(\overline{CF}+\overline{DE}\right)+2\overline{FE} = \left(\overline{CD}-\overline{FE}\right)+2\overline{FE} = \overline{CD}+\overline{FE} = a+\overline{FE},\] so the side length of square $EFGH$ is $\overline{FE} = 2b-a$, and as $F$ has coordinates $(b,0)$, $G$ therefore has coordinates $(b,2b-a)$.

 Now, by symmetry, the center of the circle is the same as the center of the square, i.e. $\left(\frac{a}{2},\frac{a}{2}\right)$, and so its radius is half of the square's diagonal, i.e. $\frac{a\sqrt{2}}{2}$. This means the equation of the circle is 
 \[\left(x-\frac{a}{2}\right)^2 + \left(y-\frac{a}{2}\right)^2 = \left(\frac{a\sqrt{2}}{2}\right)^2 = \frac{a^2}{2},\] and as $G$ lies on the circle, its coordinates must satisfy this equation, yielding 
 \[\left(b-\frac{a}{2}\right)^2 + \left(\left(2b-a\right)-\frac{a}{2}\right)^2 = \frac{a^2}{2}.\] Upon simplifying, this becomes $2a^2-7ab+5b^2 = 0$, which factors as $(2a-5b)(a-b) = 0$. Recalling that $b = \overline{DF} < \overline{DC} = a$, we cannot have $a = b$, so the solution must instead be $b = \frac{2}{5}a$. Thus the required ratio of areas is 
 \[\left(\frac{a-2b}{a}\right)^2 = \left(1-2\cdot\frac{b}{a}\right)^2 = \left(1-2\cdot\frac{2}{5}\right)^2 = \left(\frac{1}{5}\right)^2 = \frac{1}{25},\] so the answer is $10 \cdot 25 + 1 = \boxed{251}$.

 




	\item (\textit{AIME II 2001}) Let $\triangle{PQR}$ be a right triangle with $PQ = 90$, $PR = 120$, and $QR = 150$. Let $C_{1}$ be the inscribed circle. Construct $\overline{ST}$ with $S$ on $\overline{PR}$ and $T$ on $\overline{QR}$, such that $\overline{ST}$ is perpendicular to $\overline{PR}$ and tangent to $C_{1}$. Construct $\overline{UV}$ with $U$ on $\overline{PQ}$ and $V$ on $\overline{QR}$ such that $\overline{UV}$ is perpendicular to $\overline{PQ}$ and tangent to $C_{1}$. Let $C_{2}$ be the inscribed circle of $\triangle{RST}$ and $C_{3}$ the inscribed circle of $\triangle{QUV}$. The distance between the centers of $C_{2}$ and $C_{3}$ can be written as $\sqrt {10n}$. What is $n$?


 
 \textbf{Answer: }725



\textbf{Solution 1 (analytic)}


                    \begin{figure}[H]
                        \centering
                        \includegraphics[width=0.5\textwidth]{Solution/AIME_II_2001_7_Solution_Figure_1.png}
                    \end{figure}
                    Let $P = (0,0)$ be at the origin. Using the formula $A = rs$ on $\triangle PQR$, where $r_{1}$ is the \href{https://artofproblemsolving.com/wiki/index.php?title=Inradius}{inradius} (similarly define $r_2, r_3$ to be the radii of $C_2, C_3$), $s = \frac{PQ + QR + RP}{2} = 180$ is the \href{https://artofproblemsolving.com/wiki/index.php?title=Semiperimeter}{semiperimeter}, and $A = \frac 12 bh = 5400$ is the area, we find $r_{1} = \frac As = 30$. Or, the inradius could be directly by using the formula $\frac{a+b-c}{2}$, where $a$ and $b$ are the legs of the right triangle and $c$ is the hypotenuse. (This formula should be used \textit{only for right triangles}.) Thus $ST, UV$ lie respectively on the lines $y = 60, x = 60$, and so $RS = 60, UQ = 30$.  

 Note that $\triangle PQR \sim \triangle STR \sim \triangle UQV$. Since the ratio of corresponding lengths of similar figures are the same, we have 

 
 \[\frac{r_{1}}{PR} = \frac{r_{2}}{RS} \Longrightarrow r_{2} = 15\ \text{and} \ \frac{r_{1}}{PQ} = \frac{r_{3}}{UQ} \Longrightarrow r_{3} = 10.\]

 Let the centers of $\odot C_2, C_3$ be $O_2 = (0 + r_{2}, 60 + r_{2}) = (15, 75), O_3 = (60 + r_{3}, 0 + r_{3}) = (70,10)$, respectively; then by the \href{https://artofproblemsolving.com/wiki/index.php?title=Distance\_formula}{distance formula} we have $O_2O_3 = \sqrt{55^2 + 65^2} = \sqrt{10 \cdot 725}$. Therefore, the answer is $n = \boxed{725}$.

 





\textbf{Solution 2 (synthetic)}


                    \begin{figure}[H]
                        \centering
                        \includegraphics[width=0.5\textwidth]{Solution/AIME_II_2001_7_Solution_Figure_2.png}
                    \end{figure}
                    We compute $r_1 = 30, r_2 = 15, r_3 = 10$ as above. Let $A_1, A_2, A_3$ respectively the points of tangency of $C_1, C_2, C_3$ with $QR$. 

 By the \href{https://artofproblemsolving.com/wiki/index.php?title=Two\_Tangent\_Theorem}{Two Tangent Theorem}, we find that $A_{1}Q = 60$, $A_{1}R = 90$. Using the similar triangles, $RA_{2} = 45$, $QA_{3} = 20$, so $A_{2}A_{3} = QR - RA_2 - QA_3 = 85$. Thus $(O_{2}O_{3})^{2} = (15 - 10)^{2} + (85)^{2} = 7250\implies n=\boxed{725}$.

 


 





\textbf{Solution 3}

The radius of an incircle is $r=A_t/\text{semiperimeter}$. The area of the triangle is equal to $\frac{90\times120}{2} = 5400$ and the semiperimeter is equal to $\frac{90+120+150}{2} = 180$. The radius, therefore, is equal to $\frac{5400}{180} = 30$. Thus using similar triangles the dimensions of the triangle circumscribing the circle with center $C_2$ are equal to $120-2(30) = 60$, $\frac{1}{2}(90) = 45$, and $\frac{1}{2}\times150 = 75$. The radius of the circle inscribed in this triangle with dimensions $45\times60\times75$ is found using the formula mentioned at the very beginning. The radius of the incircle is equal to $15$. 

 
Defining $P$ as $(0,0)$, $C_2$ is equal to $(60+15,15)$ or $(75,15)$. Also using similar triangles, the dimensions of the triangle circumscribing the circle with center $C_3$ are equal to $90-2(30)$, $\frac{1}{3}\times120$, $\frac{1}{3}\times150$ or $30,40,50$. The radius of $C_3$ by using the formula mentioned at the beginning  is $10$. Using $P$ as $(0,0)$, $C_3$ is equal to $(10, 60+10)$ or $(10,70)$. Using the distance formula, the distance between $C_2$ and $C_3$: $\sqrt{(75-10)^2 +(15-70)^2}$ this equals $\sqrt{7250}$ or $\sqrt{725\times10}$, thus $n$ is $\boxed{725}$.

 \textbf{Note} The problem can be reduced to a $3-4-5$ triangle for the initial calculations, where $C_2$ is calculated to be ($\frac{5}{2}, \frac{1}{2})$, and $C_3$ is calculated to be ($\frac{1}{3}, \frac{7}{3})$. After we find the incenters the points can be scaled up by a factor of $30$ for the final distance calculation.

 





\textbf{Solution 4 (easy but hard to see)}

We can calculate the inradius of each triangle as with the previous solutions. Now, notice that the pentagon $PSTVU$ is a square with its corner cut off. We literally complete the square and mark the last corner as point X. Now, construct right triangle with $C_3C_2$ as its hypotenuse. The right angle will be at point $Y$. We will now find the length of each leg and use the Pythagorean Theorem to compute the desired length. We see that the length of $C_3Y$ is $50 + 15 = 65$, as seen by the inradius of $C_2$ and $10$ less than the square's side length. $C_2Y$ is $45 + 10 = 55$, which is $15$ less than the square plus the inradius of $C_3$. Our final answer is $\sqrt{65^2 + 55^2} = \sqrt{7250} = \sqrt{(10)(725)} \rightarrow \boxed{725}$.

 




	\item (\textit{AIME II 2001}) A certain function $f$ has the properties that $f(3x) = 3f(x)$ for all positive real values of $x$, and that $f(x) = 1 - |x - 2|$ for $1\leq x \leq 3$. Find the smallest $x$ for which $f(x) = f(2001)$.


 
 \textbf{Answer: }429



\textbf{Solution 1}

Iterating the condition $f(3x) = 3f(x)$, we find that $f(x) = 3^kf\left(\frac{x}{3^k}\right)$ for positive integers $k$. We know the definition of $f(x)$ from $1 \le x \le 3$, so we would like to express $f(2001) = 3^kf\left(\frac{2001}{3^k}\right),\ 1 \le \frac{2001}{3^k} \le 3 \Longrightarrow k = 6$. Indeed,

 
 \[f(2001) = 729\left[1 - \left| \frac{2001}{729} - 2\right|\right] = 186.\]

 We now need the smallest $x$ such that $f(x) = 3^kf\left(\frac{x}{3^k}\right) = 186$. The \href{https://artofproblemsolving.com/wiki/index.php?title=Range}{range} of $f(x),\ 1 \le x \le 3$, is $0 \le f(x) \le 1$. So when $1 \le \frac{x}{3^k} \le 3$, we have $0 \le f\left(\frac{x}{3^k}\right) = \frac{186}{3^k} \le 1$. Multiplying by $3^k$: $0 \le 186 \le 3^k$, so the smallest value of $k$ is $k = 5$. Then,

 
 \[186 = {3^5}f\left(\frac{x}{3^5}\right).\]

 Because we forced $1 \le \frac{x}{3^5} \le 3$, so

 
 \[186 = {3^5}f\left(\frac{x}{3^5}\right) = 243\left[1 - \left| \frac{x}{243} - 2\right|\right] \Longrightarrow x = \pm 57 + 2  \cdot 243.\]

 We want the smaller value of $x = \boxed{429}$.

 An alternative approach is to consider the graph of $f(x)$, which iterates every power of $3$, and resembles the section from $1 \le x \le 3$ dilated by a factor of $3$ at each iteration.

 





\textbf{Solution 2 (Graphing)}


                    \begin{figure}[H]
                        \centering
                        \includegraphics[width=0.5\textwidth]{Solution/AIME_II_2001_8_Solution_Figure_1.png}
                    \end{figure}
                    First, we start by graphing the function when $1\leq{x}\leq3$, which consists of the lines $y=x-1$ and $y=3-x$ that intersect at $(2,1)$. Similarly, using $f(3x)=3f(x)$, we get a dilation of our initial figure by a factor of 3 for the next interval and so on.Observe that the intersection of two lines always has coordinates $(2y,y)$ where $y=3^a$ for some $a$. First, we compute $f(2001)$. The nearest intersection point is $(1458,729)$ when $a=7$. Therefore, we can safely assume that $f(2001)$ is somewhere on the line with a slope of $-1$ that intersects at that nearest point. Using the fact that the slope of the line is $-1$, we compute $f(2001)=729-543=186$. However, we want the minimum value such that $f(x)=186$ and we see that there is another intersection point on the left which has a $y>186$, namely $(486,243)$. Therefore, we want the point that lies on the line with slope $1$ that intersects this point. Once again, since the slope of the line is $1$, we get $x=486-57=\boxed{429}$.

 





\textbf{Solution 3 (Complete Bash but FAST)}

We evaluate the first few terms of f(x) to try to find a pattern. 

 F(1)=0F(2)=1F(3)=0F(4)=1F(5) = 3(F($\frac{5}{3}$)) = 2

 That doesn‘t seem to be getting us anywhere. We notice what we did with f(5) will probably work with f(2001). 

 $F(2001) = 3f(667)=9f(\frac{667}{3}) = 27f(\frac{667}{9}) = 81f(\frac{667}{27})=243f(\frac{667}{81})=729f(\frac{667}{243})$

 From here, we can evaluate f(2001) = $186$ when we plug in $\frac{667}{243}$ into $1 - |x - 2|$. So all we need to find is the least number, let‘s call it, say y such that f(y)=186. 

 Repeating the same process we did before with f(2001), 

 $186 = F(y)= 3f(\frac{y}{3}) = 9f(\frac{y}{9}) = 27f(\frac{y}{27})=81f(\frac{y}{81}) = 243f(\frac{y}{243})$

 Notice that we stopped at $243f(\frac{y}{243})$ because $\frac{186}{243}$ is inside the range of $1-|x-2|$, which is [0,1]. Now, f(y/243) = 186/243. Setting $186/243 = 1-|x-2|$, we get 2 solutions for x:  $\frac{543}{243}$ and $\frac{429}{243}$. 

 Now, the problem asks for the smallest solution, so we obviously choose 429/243 as the solution for 243f(y/243) because it is smaller. 

 We found that $\frac{y}{243}=\frac{429}{243}$, and solving this equation gives our answer $\boxed{429}$

 by Coach Jay

 




	\item (\textit{AIME II 2001}) Each unit square of a 3-by-3 unit-square grid is to be colored either blue or red. For each square, either color is equally likely to be used. The probability of obtaining a grid that does not have a 2-by-2 red square is $\frac {m}{n}$, where $m$ and $n$ are relatively prime positive integers. Find $m + n$.


 
 \textbf{Answer: }929



\textbf{Solution 1}

We can use \href{https://artofproblemsolving.com/wiki/index.php?title=Complementary\_counting}{complementary counting}, counting all of the colorings that have at least one red $2\times 2$ square.

 
                \begin{itemize}
                    \item For at least one red $2 \times 2$ square:
                \end{itemize}


                
            \begin{enumerate}[label=(\roman*)]
                \item There are four $2 \times 2$ squares to choose which one will be red. Then there are $2^5$ ways to color the rest of the squares. $4 \cdot 32=128$
            \end{enumerate}
            
                \begin{itemize}
                    \item For at least two $2 \times 2$ squares:
                \end{itemize}


                
            \begin{enumerate}[label=(\roman*)]
                \item There are two cases: those with two $2 \times 2$ red squares on one side and those without red squares on one side but at two corners.\item The first case is easy: 4 ways to choose which the side the squares will be on, and $2^3$ ways to color the rest of the squares, so 32 ways to do that. For the second case, there will be only two ways to pick two squares, and $2^2$ ways to color the other squares. $32+8=40$
            \end{enumerate}
            
                \begin{itemize}
                    \item For at least three $2 \times 2$ squares:
                \end{itemize}


                
            \begin{enumerate}[label=(\roman*)]
                \item Choosing three such squares leaves only one square left, with four places to place it. This is $2 \cdot 4 = 8$ ways.
            \end{enumerate}
            
                \begin{itemize}
                    \item For at least four $2 \times 2$ squares, we clearly only have one way.
                \end{itemize}


                By the \href{https://artofproblemsolving.com/wiki/index.php?title=Principle\_of\_Inclusion-Exclusion}{Principle of Inclusion-Exclusion}, there are (alternatively subtracting and adding) $128-40+8-1=95$ ways to have at least one red $2 \times 2$ square.

 There are $2^9=512$ ways to paint the $3 \times 3$ square with no restrictions, so there are $512-95=417$ ways to paint the square with the restriction. Therefore, the probability of obtaining a grid that does not have a $2 \times 2$ red square is $\frac{417}{512}$, and $417+512=\boxed{929}$.

 





\textbf{Solution 2}

We consider how many ways we can colour the $2 \times 2$ grid:

 $(1)$: All the grids are red--$1$ case

 $(2)$: One unit square is blue--The blue lies on the center of the bigger square, makes no $2 \times 2$ grid $9-1=8$ cases

 $(3)$: Two unit squares are blue--one of the squares lies in the center of the bigger square, makes no $2 \times 2$ grid, $8$ cases. Or, two squares lie on second column, first row, second column third row; second row first column, second row third column, 2 extra cases. $\binom 9 2-8-2=26$ cases

 $(4)$ Three unit squares are blue. We find that if a $2 \times 2$ square is formed, there are 5 extra unit squares can be painted. But cases that three squares in the same column or same row is overcomunted. So in this case, there are $4\cdot (\binom 5 3)-4=36$

 $(5)$ Four unit squares are blue, no overcounted case will be considered. there are $4\cdot \binom 5 4=20$

 $(6)$ Five unit squares are blue, $4$ cases in all

 Sum up those cases, there are $1+8+26+36+20+4=95$ cases that a $2 \times 2$ grid can be formed.

 In all, there are $2^9=512$ possible ways to paint the big square, so the answer is $1-\frac{95}{512}=\frac{417}{512}$ leads to $\boxed{929}$

 





\textbf{Solution 3 (Case Work)}


 \[\begin{array}{|c|c|c|} \hline C_{11} & C_{12} & C_{13}\\ \hline C_{21} & C_{22} & C_{23}\\ \hline C_{31} & C_{32} & C_{33}\\ \hline \end{array}\]

 Case 1: The 3-by-3 unit-square grid has exactly $1$ 2-by-2 red square

 Assume the 2-by-2 red square is at $C_{11}, C_{12}, C_{21}, C_{22}$. To make sure there are no more 2-by-2 red squares, $C_{31} \text{and} C_{32}$ can't both be red and $C_{13} \text{and} C_{23}$ can't both be red. Meaning that there are $2^2-1=3$ coloring methods for $C_{31} \text{and} C_{32}$ and $C_{13} \text{and} C_{23}$. $C_{33}$ can be colored with either colors. However, the coloring method where $C_{23}, C_{32}, C_{33}$ are all red needs to be removed. For exactly one 2-by-2 red square at $C_{11}, C_{12}, C_{21}, C_{22}$, there are $3 \cdot 3 \cdot 2 -1=17$ coloring methods. As there are $4$ locations for the 2-by-2 red square on the 3-by-3 unit-square grid, there are $17 \cdot 4 = 68$ coloring methods.

 
Case 2: The 3-by-3 unit-square grid has exactly $2$ 2-by-2 red squares

 Case 2.1: $2$ 2-by-2 red squares take up $6$ unit grids

 Assume the $2$ 2-by-2 red squares are at $C_{11}, C_{12}, C_{21}, C_{22}, C_{31}, C_{32}$. To make sure there are no more 2-by-2 red squares, $C_{13} \text{and} C_{23}$ can't both be red. Meaning that there are $2^2-1=3$ coloring methods for $C_{13} \text{and} C_{23}$. $C_{33}$ can be colored with either colors. However, the coloring method where $C_{13}, C_{23}, C_{33}$ are all red needs to be removed. For exactly $2$ 2-by-2 red squares at $C_{11}, C_{12}, C_{21}, C_{22}, C_{31}, C_{32}$, there are $3 \cdot 2 -1=5$ coloring methods. As there are $4$ locations for the $2$ 2-by-2 red squares on the 3-by-3 unit-square grid, there are $5 \cdot 4 = 20$ coloring methods.

 Case 2.2: $2$ 2-by-2 red squares take up $7$ unit grids

 Assume the $2$ 2-by-2 red squares are at $C_{11}, C_{12}, C_{21}, C_{22}, C_{23}, C_{32}, C_{33}$. To make sure there are no more 2-by-2 red squares, $C_{13}$ and $C_{31}$ can't be red. Meaning that there is $1$ coloring method for $C_{13}$ and $C_{31}$. For exactly $2$ 2-by-2 red squares at $C_{11}, C_{12}, C_{21}, C_{22}, C_{23}, C_{32}, C_{33}$, there is $1$ coloring method. As there are $2$ locations for the $2$ 2-by-2 red squares on the 3-by-3 unit-square grid, there are $1 \cdot 2 = 2$ coloring methods.

 Hence, if the 3-by-3 unit-square grid has exactly $2$ 2-by-2 red squares, there are $20+2 = 22$ coloring methods.

 
Case 3: The 3-by-3 unit-square grid has exactly $3$ 2-by-2 red squares

 Assume the $3$ 2-by-2 red squares are at $C_{11}, C_{12}, C_{13}, C_{21}, C_{22}, C_{23}, C_{32}, C_{33}$. To make sure there are no more 2-by-2 red squares, $C_{33}$ can't be red. Meaning that there is $1$ coloring method for $C_{33}$. For exactly $3$ 2-by-2 red squares at $C_{11}, C_{12}, C_{13}, C_{21}, C_{22}, C_{23}, C_{32}, C_{33}$, there is $1$ coloring method. As there are $4$ locations for the $3$ 2-by-2 red squares on the 3-by-3 unit-square grid, there are $1 \cdot 4 = 4$ coloring methods.

 
Case 4: The 3-by-3 unit-square grid has exactly $4$ 2-by-2 red squares

 If the 3-by-3 unit-square grid has exactly $4$ 2-by-2 red squares, all $9$ unit grids are red and there is $1$ coloring method.

 
In total, there are $68+22+4+1=95$ coloring methods with 2-by-2 red squares.
 \[\frac{m}{n}=1-\frac{95}{2^9}=\frac{417}{512}\]

 
 \[m+n=417+512=\boxed{\textbf{929}}\]

 




	\item (\textit{AIME II 2001}) How many positive integer multiples of 1001 can be expressed in the form $10^{j} - 10^{i}$, where $i$ and $j$ are integers and $0\leq i < j \leq 99$?


 
 \textbf{Answer: }784



\textbf{Solution 1}

The \href{https://artofproblemsolving.com/wiki/index.php?title=Prime\_factorization}{prime factorization} of $1001 = 7\times 11\times 13$. We have $7\times 11\times 13\times k = 10^j - 10^i = 10^i(10^{j - i} - 1)$. Since $\text{gcd}\,(10^i = 2^i \times 5^i, 7 \times 11 \times 13) = 1$, we require that $1001 = 10^3 + 1 | 10^{j-i} - 1$. From the factorization $10^6 - 1 = (10^3 + 1)(10^{3} - 1)$, we see that $j-i = 6$ works; also, $a-b | a^n - b^n$ implies that $10^{6} - 1 | 10^{6k} - 1$, and so any $\boxed{j-i \equiv 0 \pmod{6}}$ will work. 

 To show that no other possibilities work, suppose $j-i \equiv a \pmod{6},\ 1 \le a \le 5$, and let $j-i-a = 6k$. Then we can write $10^{j-i} - 1 = 10^{a} \left(10^{6k} - 1\right) + \left(10^{a} - 1\right)$, and we observe that $10^{6k}-1 = \left(10^6\right)^k-1^k$, which is divisible by $10^6-1$, and thus by $10^3+1$ (as $10^6-1 = \left(10^3\right)^2-1^2 = \left(10^3+1\right)\left(10^3-1\right)$). It follows that the first term is divisible by $1001$, whereas we can easily verify that the second term, $10^a-1$, is not divisible by $1001$ for $1 \le a \le 5$, so $10^{j-i}$ is not divisible by $1001$ for such $a$.

 If $j - i = 6, j\leq 99$, then we can have solutions of $10^6 - 10^0, 10^7 - 10^1, \dots\implies 94$ ways. If $j - i = 12$, we can have the solutions of $10^{12} - 10^{0},\dots\implies 94 - 6 = 88$, and so forth. Therefore, the answer is $94 + 88 + 82 + \dots + 4\implies 16\left(\dfrac{98}{2}\right) = \boxed{784}$.

 





\textbf{Solution 2}

Observation: We see that there is a pattern with $10^k \pmod{1001}$. 
 \[10^0 \equiv 1 \pmod{1001}\]
 \[10^1 \equiv 10 \pmod{1001}\]
 \[10^2 \equiv 100 \pmod{1001}\]
 \[10^3 \equiv -1 \pmod{1001}\]
 \[10^4 \equiv -10 \pmod{1001}\]
 \[10^5 \equiv -100 \pmod{1001}\]
 \[10^6 \equiv 1 \pmod{1001}\]
 \[10^7 \equiv 10 \pmod{1001}\]
 \[10^8 \equiv 100 \pmod{1001}\]

 So, this pattern repeats every 6. 

 Also, $10^j-10^i \equiv 0 \pmod{1001}$, so$10^j \equiv 10^i \pmod{1001}$, and thus, 
 \[j \equiv i \pmod{6}\]. Continue with the 2nd paragraph of solution 1, and we get the answer of $\boxed{784}$.

 -AlexLikeMath

 





\textbf{Solution 3}

Note that $1001=7\cdot 11\cdot 13,$ and note that $10^6 \equiv 1 \pmod{p}$ for prime $p \mid 1001$; therefore, the order of 10 modulo $7,11$, and $13$ must divide 6. A quick check on 7 reveals that it is indeed 6. Therefore we note that $i-j=6k$ for some natural number k. From here, we note that for $j=0,1,2,3,$ we have 16 options and we have 15,14,...,1 option(s) for the next 90 numbers (6 each), so our total is $4\cdot 16 + 6 \cdot \frac{15 \cdot 16}{2} = \boxed{784}$.

 





\textbf{Solution 4}

$10^j - 10^i \equiv 0 \pmod{1001} \iff 10^{j - i} - 1 \equiv 0 \pmod{1001} \iff 10^{j - i} \equiv 1 \pmod{1001} \iff j \equiv i \pmod 6$, by the same argument as in all the solutions above. If $j \equiv i \equiv n \pmod 6$ for $n = 0, 1, 2, 3$, there are $17$ choices for each value of $n$, yielding $4 \cdot \dbinom{17}{2} = 544$. However, if $n = 4, 5$, there are only $16$ choices, giving us $2 \cdot \dbinom{16}{2} = 240$. So, our final answer is $544 + 240 = \boxed{784}$.~Puck\_0

 




	\item (\textit{AIME II 2001}) Club Truncator is in a soccer league with six other teams, each of which it plays once. In any of its 6 matches, the probabilities that Club Truncator will win, lose, or tie are each $\frac {1}{3}$. The probability that Club Truncator will finish the season with more wins than losses is $\frac {m}{n}$, where $m$ and $n$ are relatively prime positive integers. Find $m + n$.


 
 \textbf{Answer: }341



\textbf{Solution 1}

Note that the probability that Club Truncator will have more wins than losses is equal to the probability that it will have more losses than wins; the only other possibility is that they have the same number of wins and losses. Thus, by the \href{https://artofproblemsolving.com/wiki/index.php?title=Complement\_principle}{complement principle}, the desired probability is half the probability that Club Truncator does not have the same number of wins and losses.

 The possible ways to achieve the same number of wins and losses are $0$ ties, $3$ wins and $3$ losses; $2$ ties, $2$ wins, and $2$ losses; $4$ ties, $1$ win, and $1$ loss; or $6$ ties. Since there are $6$ games, there are $\frac{6!}{3!3!}$ ways for the first, and $\frac{6!}{2!2!2!}$, $\frac{6!}{4!}$, and $1$ ways for the rest, respectively, out of a total of $3^6$. This gives a probability of $141/729$. Then the desired answer is $\frac{1 - \frac{141}{729}}{2} = \frac{98}{243}$, so the answer is $m+n = \boxed{341}$.

 





\textbf{Solution 2}

At first, it wins $6$ games, only one way.

 Secondly, it wins $5$ games, the other game can be either win or loss, there are $\binom{6}{5}\cdot 2=12$ ways.

 Thirdly, it wins $4$ games, still the other two games can be either win or loss, there are $\binom{6}{4}\cdot 2^2=60$ ways.

 Fourthly, it wins $3$ games, this time, it can't lose $3$ games but other arrangements of the three non-winning games are fine, there are $\binom{6}{3}\cdot (2^3-1)=140$ ways.

 Fifth case, it wins $2$ games, only $0/1$ lose and  $4/3$ draw is ok, so there are $\binom{6}{2}(1+\binom{4}{1})=75$ cases.

 Last case, it only wins $1$ game so the rest games must be all draw, $1$ game.

 The answer is $\frac{1+12+60+140+75+6}{3^6}=\frac{98}{243}$ leads to $\boxed{341}$

 





\textbf{Solution 3}

This problem is symmetrical, and we should probably take advantage of that. We have that:

 
 \begin{align*} 1 &= P(\#\text{ of wins} > \text{losses}) + P(\#\text{ of wins} < \text{losses}) + P(\#\text{ of wins} = \text{losses}) \\   &= 2P(\#\text{ of wins} > \text{losses}) + P(\#\text{ of wins} = \text{losses}) \end{align*}


 Now, all we need to do is find the probability that the number of wins is equal to the number of losses, and then it's simple linear equations from there! We have 4 cases. We can let $W/L/T$ be the format for these cases. For example, $3/2/1$ would mean 3 wins, 2 losses, and 1 tie. Now, we begin our casework.

 
 \[C_1: 3/3/0 \Longrightarrow \binom{6}{3} \cdot \left(\frac{1}{3}\right)^6 = \frac{20}{729}\]

 
 \[C_2: 2/2/2 \Longrightarrow \binom{6}{2}\binom{4}{2} \cdot \left(\frac{1}{3}\right)^6 = \frac{90}{729}\]

 
 \[C_3: 1/1/4 \Longrightarrow \binom{6}{4} \cdot 2 \cdot \left(\frac{1}{3}\right)^6 = \frac{30}{729}\]

 
 \[C_4: 0/0/6 \Longrightarrow \left(\frac{1}{3}\right)^6 = \frac{1}{729}\]

 Summing these up we get $\dfrac{141}{729} = \dfrac{47}{243}$. Now, we have $2x + \dfrac{47}{243} = 1$, where $x = P(\#\text{ of wins} > \text{losses})$. Solving, we get $x = \dfrac{98}{243} \Longrightarrow \boxed{341}$

 -jb2015007

 




	\item (\textit{AIME II 2001}) Given a triangle, its midpoint triangle is obtained by joining the midpoints of its sides. A sequence of polyhedra $P_{i}$ is defined recursively as follows: $P_{0}$ is a regular tetrahedron whose volume is 1. To obtain $P_{i + 1}$, replace the midpoint triangle of every face of $P_{i}$ by an outward-pointing regular tetrahedron that has the midpoint triangle as a face. The volume of $P_{3}$ is $\frac {m}{n}$, where $m$ and $n$ are relatively prime positive integers. Find $m + n$.


 
 \textbf{Answer: }101



\textbf{Solution}

On the first construction, $P_1$, four new tetrahedra will be constructed with side lengths $\frac 12$ of the original one. Since the ratio of the volume of similar polygons is the cube of the ratio of their corresponding lengths, it follows that each of these new tetrahedra will have volume $\left(\frac 12\right)^3 = \frac 18$. The total volume added here is then $\Delta P_1 = 4 \cdot \frac 18 = \frac 12$.

 We now note that for each midpoint triangle we construct in step $P_{i}$, there are now $6$ places to construct new midpoint triangles for step $P_{i+1}$. The outward tetrahedron for the midpoint triangle provides $3$ of the faces, while the three equilateral triangles surrounding the midpoint triangle provide the other $3$. This is because if you read this question carefully, it asks to add new tetrahedra to each face of $P_{i}$ which also includes the ones that were left over when we did the previous addition of tetrahedra. However, the volume of the tetrahedra being constructed decrease by a factor of $\frac 18$. Thus we have the recursion $\Delta P_{i+1} = \frac{6}{8} \Delta P_i$, and so $\Delta P_i = \frac 12 \cdot \left(\frac{3}{4}\right)^{i-1} P_1$.

 The volume of $P_3 = P_0 + \Delta P_1 + \Delta P_2 + \Delta P_3 = 1 + \frac 12 + \frac 38 + \frac 9{32} = \frac{69}{32}$, and $m+n=\boxed{101}$. Note that the summation was in fact a \href{https://artofproblemsolving.com/wiki/index.php?title=Geometric\_series}{geometric series}.

 




	\item (\textit{AIME II 2001}) In quadrilateral $ABCD$, $\angle{BAD}\cong\angle{ADC}$ and $\angle{ABD}\cong\angle{BCD}$, $AB = 8$, $BD = 10$, and $BC = 6$. The length $CD$ may be written in the form $\frac {m}{n}$, where $m$ and $n$ are relatively prime positive integers. Find $m + n$.


 
 \textbf{Answer: }069



\textbf{Solution 1}

Notice that $\angle BAD = \angle ADC$ and $\angle ABD = \angle DCB$, which motivates us to find similar triangles. We extend $CB$ and $DA$ to point $E$. We know that $\triangle ABD \sim \triangle DCE$. Hence $\angle ADB = \angle DEC$, and $\triangle BDE$ is \href{https://artofproblemsolving.com/wiki/index.php?title=Isosceles\_triangle}{isosceles}. Then $BD = BE = 10$. 

 
                    \begin{figure}[H]
                        \centering
                        \includegraphics[width=0.5\textwidth]{Solution/AIME_II_2001_13_Solution_Figure_1.png}
                    \end{figure}
                    Using the similarity, we have:

 
 \[\frac{AB}{BD} = \frac 8{10} = \frac{CD}{CE} = \frac{CD}{16} \Longrightarrow CD = \frac{64}5\]

 The answer is $m+n = \boxed{069}$.

 
\textbf{Extension}: To Find $AD$, use Law of Cosines on $\triangle BCD$ to get $\cos(\angle BCD)=\frac{13}{20}$Then since $\angle BCD=\angle ABD$ use Law of Cosines on $\triangle ABD$ to find $AD=2\sqrt{15}$

 





\textbf{Solution 2}

Draw a line from $B$, parallel to $\overline{AD}$, and let it meet $\overline{CD}$ at $M$. Note that $\triangle{DAB}$ is similar to $\triangle{BMC}$ by AA similarity, since $\angle{ABD}=\angle{MCB}$ and since $BM$ is parallel to $CD$ then $\angle{BMC}=\angle{ADM}=\angle{DAB}$. Now since $ADMB$ is an isosceles trapezoid, $MD=8$. By the similarity, we have $MC=AB\cdot \frac{BC}{BD}=8\cdot \frac{6}{10}=\frac{24}{5}$, hence $CD=MC+MD=\frac{24}{5}+8=\frac{64}{5}\implies 64+5=\boxed{069}$.

 





\textbf{Solution 3}

Since $\angle{BAD}=\angle{ADM}$, if we extend AB and DC, they must meet at one point to form a isosceles triangle $\triangle{ADM}$.Now, since the problem told that $\angle{ABD}=\angle{BCD}$, we can imply that $\angle{DBM}=\angle{BCM}$Since $\angle{M}=\angle{M}$, so $\triangle{CBM}\sim\triangle{BDM}$. Assume the length of $BM=x$;Since $\frac{BC}{MB}=\frac{DB}{MD}$ we can get $\frac{6}{x}=\frac{10}{8+x}$, we get that $x=12$.So $AM=DM=20$ similarly, we use the same pair of similar triangle we get $\frac{CM}{BM}=\frac{BM}{DM}$, we get that $CM=\frac{36}{5}$. Finally, $CD=MD-MC=\frac{64}{5}\implies 64+5=69=\boxed{069}$~bluesoul

 





\textbf{Solution 3}

Denote $\angle{BAD}=\angle{CDA}=x$, and $\angle{ABD}=\angle{BCD}=y$. Note that $\angle{ADB}=180^\circ-x-y$, and $\angle{DBC}=360^\circ-2x-2y$. This motivates us to draw the angle bisector of $\angle{DBC}$ because $\angle{DBC} = 2 \angle{ADB}$, so we do so and consider the intersection with $CD$ as $E$. By the angle bisector theorem, we have $\frac{CE}{DE} = \frac{BC}{BD} = \frac{3}{5}$, so we write $CE=3z$ and $DE=5z$. We also know that $\angle{EBC}=\angle{ADB}$ and $\angle{BCE}=\angle{DBA}$, so $\triangle{ADB} \sim \triangle{EBC}$. Hence, $\frac{CE}{BC}=\frac{AB}{BD}$, so we have $3z=\frac{24}{5}$. As $CD=8z$, it must be that $CD=\frac{64}{5}$, so the final answer is $\boxed{069}$.

 




	\item (\textit{AIME II 2001}) There are $2n$ complex numbers that satisfy both $z^{28} - z^{8} - 1 = 0$ and $|z| = 1$. These numbers have the form $z_{m} = \cos\theta_{m} + i\sin\theta_{m}$, where $0\leq\theta_{1} < \theta_{2} < \ldots < \theta_{2n} < 360$ and angles are measured in degrees. Find the value of $\theta_{2} + \theta_{4} + \ldots + \theta_{2n}$.


 
 \textbf{Answer: }840



\textbf{Solution 1 (Solution 2 more in depth)}

If $|z| = 1$, then it is evident that both $|z^{28}| = |z^8| = 1$. Then, $z^{28} - z^8 = 1$. Say $z^{28} = a+bi$ and $z^8 = r+si$. Then, $a-r + (b-s)i = 1$, in which through equivalation of real and imaginary parts, $a=r+1$ and $b=s$, in which we have $z^{28} = (1+r)+si$ and $z^8 = r+si$. Then, because $|z^{28}|=|z^8|$, we have $(1+r)^2+s^2 = r^2+s^2$, in which $r=-\frac{1}{2}$. Then, substituting into the definition of $|z^8|$, we see that $s=\pm \frac{\sqrt{3}}{2}$. Thus, we have two cases.

 $\textbf{Case 1:}$ $z^{28}=\frac{1}{2} + \frac{\sqrt{3}}{2}i, z^8=\frac{-1}{2}+\frac{\sqrt{3}}{2}i$

 In this case, we rewrite everything in terms of $\text{cis}(x)$. Then, $z^{28} = \text{cis}(\frac{\pi}{3})$, $z^8=\text{cis}(\frac{2\pi}{3})$. Then, because each $z_m = \text{cis}(\theta_m)$, we have that $\text{cis}(28\theta)=\text{cis}(\frac{\pi}{3})$, and $\text{cis}(8\theta)=\text{cis}(\frac{2\pi}{3})$. Thus $28\theta \equiv \frac{\pi}{3} \mod 2\pi$, and $8\theta \equiv \frac{2\pi}{3} \mod 2\pi$.

 The second congruence gives us $\theta \equiv \frac{\pi}{12} \mod \frac{\pi}{4}$, or $\theta = \frac{\pi}{12} + \frac{\pi}{4}k$. Substitution into the first modulo gives us $\frac{7\pi}{3} + 7\pi k \equiv \frac{\pi}{3} \mod 2\pi$, in which subtracting $\frac{\pi}{3}$ gives us $2\pi+7\pi k\equiv 0 \mod 2\pi$. This only holds if $k=2m$, and thus we have $\theta = \frac{\pi}{12} + \frac{\pi m}{2}$. The solutions in incremental order to this case are $\frac{\pi}{12},\; \frac{7\pi}{12},\; \frac{13\pi}{12},\; \frac{19\pi}{12}$.

 $\textbf{Case 2:}$ $z^{28}=\frac{1}{2} - \frac{\sqrt{3}}{2}i, z^8=\frac{-1}{2}-\frac{\sqrt{3}}{2}i$

 Once again, conversion into $\text{cis}(x)$ form and using De Moivre's gives us $28\theta \equiv -\frac{\pi}{3} \mod 2\pi$ and $8\theta \equiv -\frac{2\pi}{3} \mod 2\pi$. Once again, the second equation gives us $\theta = -\frac{\pi}{12} + \frac{\pi}{4}k$. Substitution of $\theta$ into the first again gives us $28x=28(-\frac{\pi}{12} + \frac{\pi k}{4})=-\frac{7\pi}{3}+7\pi k$, in which reducing modulo $2\pi$ and adding $\frac{\pi}{3}$, we obtain $-\frac{2\pi}{3}+7\pi k \equiv 0 \pmod{2\pi}$. Once again, $k=2m$, and so $\theta = -\frac{\pi}{12} + \frac{\pi m}{2}$. The solutions for $\theta$ are $\frac{5\pi}{12},\;\frac{11\pi}{12},\; \frac{17\pi}{12},\; \frac{23\pi}{12}$.

 Converting everything into degrees and putting in incremental order, we see that each $\theta_m : 1\le m\le 8 = 15^\circ,\;75^\circ,\;105^\circ,\;165^\circ,\;195^\circ,\;255^\circ,\;285^\circ,\;345^\circ$.

 Adding each $\theta_{2k} : 1 \le k \le 4$, we obtain a final answer of $75+165+255+345 = \boxed{840}$.

 





\textbf{Solution 2}

$z$ can be written in the form $\text{cis\,}\theta$. Rearranging, we find that  $\text{cis\,}{28}\theta = \text{cis\,}{8}\theta+1$

 Since the real part of $\text{cis\,}{28}\theta$ is one more than the real part of $\text{cis\,} {8}\theta$ and their imaginary parts are equal, it is clear that either $\text{cis\,}{28}\theta = \frac{1}{2}+\frac {\sqrt{3}}{2}i$ and $\text{cis\,} {8}\theta =  -\frac{1}{2}+\frac {\sqrt{3}}{2}i$, or $\text{cis\,}{28}\theta = \frac{1}{2} - \frac{\sqrt{3}}{2}i$ and  $\text{cis\,} {8}\theta =  -\frac{1}{2}- \frac{\sqrt{3}}{2}i$

 
                \begin{itemize}
                    \item Case 1  : $\text{cis\,}{28}\theta = \frac{1}{2}+ \frac{\sqrt{3}}{2}i$ and $\text{cis\,} {8}\theta =  -\frac{1}{2}+\frac{\sqrt{3}}{2}i$
                \end{itemize}


                Setting up and solving equations, $Z^{28}= \text{cis\,}{60^\circ}$ and $Z^8= \text{cis\,}{120^\circ}$, we see that the solutions common to both equations have arguments $15^\circ , 105^\circ, 195^\circ,$ and $\ 285^\circ$. We can figure this out by adding 360 repeatedly to the number 60 to try and see if it will satisfy what we need. We realize that it does not work in the integer values.

 
                \begin{itemize}
                    \item Case 2  : $\text{cis\,}{28}\theta = \frac{1}{2} -\frac {\sqrt{3}}{2}i$ and $\text{cis\,} {8}\theta =  -\frac {1}{2} -\frac{\sqrt{3}}{2}i$
                \end{itemize}


                Again setting up equations ($Z^{28}= \text{cis\,}{300^\circ}$ and $Z^{8} = \text{cis\,}{240^\circ}$) we see that the common solutions have arguments of  $75^\circ, 165^\circ, 255^\circ,$ and $345^\circ$

 Listing all of these values, we find that $\theta_{2} + \theta_{4} + \ldots + \theta_{2n}$ is equal to $(75 + 165 + 255 + 345) ^\circ$ which is equal to $\boxed {840}$ degrees. We only want the sum of a certain number of theta, not all of it.

 




	\item (\textit{AIME II 2001}) Let $EFGH$, $EFDC$, and $EHBC$ be three adjacent square faces of a cube, for which $EC = 8$, and let $A$ be the eighth vertex of the cube. Let $I$, $J$, and $K$, be the points on $\overline{EF}$, $\overline{EH}$, and $\overline{EC}$, respectively, so that $EI = EJ = EK = 2$. A solid $S$ is obtained by drilling a tunnel through the cube. The sides of the tunnel are planes parallel to $\overline{AE}$, and containing the edges, $\overline{IJ}$, $\overline{JK}$, and $\overline{KI}$. The surface area of $S$, including the walls of the tunnel, is $m + n\sqrt {p}$, where $m$, $n$, and $p$ are positive integers and $p$ is not divisible by the square of any prime. Find $m + n + p$.


 
 \textbf{Answer: }417



\textbf{Solution}


                    \begin{figure}[H]
                        \centering
                        \includegraphics[width=0.5\textwidth]{Solution/AIME_II_2001_15_Solution_Figure_1.png}
                    \end{figure}
                                 
                    \begin{figure}[H]
                        \centering
                        \includegraphics[width=0.5\textwidth]{Solution/AIME_II_2001_15_Solution_Figure_2.png}
                    \end{figure}
                    Set the coordinate system so that vertex $E$, where the drilling starts, is at $(8,8,8)$. Using a little visualization (involving some \href{https://artofproblemsolving.com/wiki/index.php?title=Similar\_triangles}{similar triangles}, because we have parallel lines) shows that the tunnel meets the bottom face (the xy plane one) in the line segments joining $(1,0,0)$ to $(2,2,0)$, and $(0,1,0)$ to $(2,2,0)$, and similarly for the other three faces meeting at the origin (by symmetry). So one face of the tunnel is the polygon with vertices (in that order), $S(1,0,0), T(2,0,2), U(8,6,8), V(8,8,6), W(2,2,0)$, and the other two faces of the tunnel are congruent to this shape. 

 Observe that this shape is made up of two congruent \href{https://artofproblemsolving.com/wiki/index.php?title=Trapezoid}{trapezoids} each with height $\sqrt {2}$ and bases $7\sqrt {3}$ and $6\sqrt {3}$. Together they make up an area of $\sqrt {2}(7\sqrt {3} + 6\sqrt {3}) = 13\sqrt {6}$. The total area of the tunnel is then $3\cdot13\sqrt {6} = 39\sqrt {6}$. Around the corner $E$ we're missing an area of $6$, the same goes for the corner opposite $E$ . So the outside area is $6\cdot 64 - 2\cdot 6 = 372$. Thus the the total surface area is $372 + 39\sqrt {6}$, and the answer is $372 + 39 + 6 = \boxed{417}$.

 




\end{enumerate}